\section{Predicates, support and binding}
\label{sec:perspectives}

The distinction between a carrier and its action becomes particularly useful
when the carrier represents logical information. Giving truth values the
trivial action still leaves a nontrivial action on predicates. For a
permutation set $X$, conjugation on $P : X \to \mathrm{Prop}$ gives
\[
  (\pi \act P)(x) \;\longleftrightarrow\; P(\pi^{-1}\act x).
\]
For example, the atom predicate $P_a(x) \equiv (x=a)$ is sent to
$P_{\pi a}$. The singleton $\{a\}$ supports $P_a$, although exchanging $a$
with another atom changes the predicate. Pitts' nominal powerset captures
precisely this distinction: its elements are supported subsets, while
equivariant subsets have empty support~\cite[Section~2.5]{pitts2013}.
There are consequently two different roles for a predicate in reasoning
about binders. As a semantic object, it may carry a finite-support
hypothesis; as the motive of an induction theorem, it may be arbitrary.
Pitts' alpha-structural induction theorem uses supported predicate
families~\cite[Theorem~8.21]{pitts2013}. Nominal2's strong induction principle
instead generalizes each recursive hypothesis over a finitely supported
context, without requiring finite support of the motive
itself~\cite[Section~7, equation~(7.1)]{urbanKaliszyk2012}. These are different
theorem interfaces: the context specifies what binders must avoid, and the
generalized hypotheses provide the strength needed to change that context.
A calculus of supported predicates therefore need not restrict the ordinary
propositions available to an induction proof.

The treatment of support witnesses gives another useful comparison.
The constructive Rocq account of Paranhos and Ventura equips nominal
carriers with a function supplying a finite support and uses setoid equality
with explicit respectfulness laws~\cite[Section~3]{paranhosVentura2025}.
In the present formulation, finite supportedness is an existential
proposition about an element of a selected action. A proof can introduce a
convenient bound, enlarge it, transport it or combine it with another bound,
without making that choice part of the element's identity. This distinction
matters even in classical mathematics: a supplied support may contain atoms
on which the element does not depend. Over infinite atoms, the intersection
theorem turns an inclusion-minimal bound into the unique least support
(Proposition~\ref{prop:least-support}), as in the classical theory of
support~\cite[Section~2.1]{pitts2013}. This construction applies to an
individually supported element; nominality of the whole carrier only supplies
such certificates uniformly. The elementary support calculus itself requires
no operation selecting a least bound. Lean's ordinary
equality and classical foundations make extensional equality and quotient
constructions available; the Rocq formulation instead exposes the chosen
equivalence relation in its action laws. Support arguments must respect the
equality used for the carrier in either setting.

Nominal2 also illustrates the additional mathematical content needed to
connect such a semantic foundation to a language specification. Urban and
Kaliszyk construct raw syntax, define alpha-equivalence, and pass to
quotients, deriving support equations and strong induction through generated
proofs~\cite[Sections~4--7]{urbanKaliszyk2012}. This passage depends on
properties of the particular binding specification: alpha-equivalence must
be an equivariant equivalence relation, and operations descending to the
quotient must respect it. Identifying semantic support with free atoms is a
further theorem about that syntax. The selected-action support laws studied
here supply reusable premises for these arguments, independently of the
constructors chosen for a language. They do not determine which occurrences
a constructor binds. In particular, an action on terms and finite support
alone do not justify a fresh-binder induction rule: the rule must explain
how changing a representative preserves the term and supplies the required
recursive cases. Separating these obligations gives a precise connection
between permutation semantics and the familiar practice of choosing bound
names fresh for the parameters of a proof.
